\chapter{CAPITULO II - IDENTIFICACIÓN}

\section{Diagnóstico}

Huayllarcocha es una comunidad ubicada en la región Cusco, Perú, que se caracteriza por ser un asentamiento rural con características propias de las comunidades altoandinas de la región. Esta ubicación geográfica determina condiciones específicas de acceso a servicios básicos, infraestructura limitada y desafíos particulares en la prestación de servicios públicos esenciales como el agua.


La situación actual del territorio de Huayllarcocha presenta características que requieren un análisis profundo y detallado. El territorio se encuentra en una zona rural del Cusco, donde las condiciones geográficas, climáticas y de acceso condicionan significativamente la disponibilidad y calidad de los servicios de agua. Las comunidades rurales de la región del Cusco enfrentan históricamente brechas importantes en la cobertura de servicios básicos, y Huayllarcocha no es ajena a esta realidad. El análisis territorial debe considerar la ubicación específica, las características topográficas, el acceso a vías de comunicación, la distancia a centros urbanos y la disponibilidad de recursos naturales, especialmente el agua.


La población de Huayllarcocha presenta características demográficas y socioeconómicas que deben ser comprendidas en su totalidad para diseñar una intervención adecuada. Las comunidades rurales del Cusco suelen tener poblaciones con NBI (Necesidad Básica Insatisfecha) significativo, estratificación socioeconómica baja y condiciones de pobreza que limitan el acceso a servicios de calidad. La población afectada por la carencia de monitoreo del agua comprende tanto a los habitantes locales que dependen directamente del agua para sus actividades domésticas, agrícolas y de subsistencia, como a los estudiantes que podrían estar vinculados a instituciones educativas de la comunidad. Es fundamental identificar el total de la población que presenta la necesidad y/o problema que justifica el programa, incluyendo sus características socioeconómicas y demográficas específicas.

El estado de la infraestructura existente para el servicio de agua en Huayllarcocha constituye un elemento crítico del diagnóstico. En las zonas rurales del Cusco, la infraestructura de agua suele presentar condiciones de deterioro, cobertura incompleta, falta de mantenimiento y ausencia de sistemas de monitoreo adecuados. El diagnóstico debe evaluar específicamente la cobertura del servicio de agua, la continuidad del servicio, la eficiencia en la distribución y la calidad del agua suministrada. Estos cuatro elementos (cobertura, continuidad, eficiencia y calidad) son indicadores fundamentales para determinar el estado general de la prestación de los servicios públicos en la comunidad.

La infraestructura hídrica existente en Huayllarcocha probablemente incluye sistemas de captación, almacenamiento y distribución que pueden ser de tipo tradicional o convencional, pero que carecen de sistemas modernos de monitoreo. En muchas comunidades rurales del Cusco, los sistemas de agua se construyeron con tecnologías simples que no incorporan sensores, dispositivos de medición automática ni sistemas de transmisión de datos en tiempo real. Esta carencia tecnológica limita significativamente la capacidad de la comunidad para conocer el estado actual de sus recursos hídricos, detectar problemas de calidad, identificar pérdidas en el sistema o tomar decisiones basadas en datos precisos sobre el uso del agua.

El diagnóstico debe considerar también el estado de los servicios públicos existentes en relación con la cobertura, continuidad, eficiencia y calidad del servicio de agua. La cobertura se refiere al porcentaje de la población que tiene acceso al servicio de agua; la continuidad indica si el servicio se proporciona de manera constante o intermitente; la eficiencia evalúa si el sistema distribuye el agua de manera óptima sin pérdidas significativas; y la calidad determina si el agua cumple con los estándares de potabilidad y seguridad necesarios para el consumo humano y otras actividades.


La situación actual de Huayllarcocha también debe analizarse en el contexto de las condiciones físicas, económicas y sociales del municipio y del área objeto de intervención. Esto incluye considerar los usos del suelo, las condiciones sociales, la salud pública, los aspectos educativos, las organizaciones cívicas, el nivel de ingresos de la población, la disponibilidad de recursos humanos y materiales en la región, y otros factores que influyen en la capacidad de la comunidad para implementar y mantener sistemas de monitoreo del agua.


Las características socioeconómicas de la población de Huayllarcocha son fundamentales para entender las limitaciones que enfrenta la comunidad. La población actual, su estratificación socioeconómica, el índice de NBI, la población en situación de miseria, y otros indicadores de bienestar social deben ser documentados en el diagnóstico. Estas características determinan la capacidad de la comunidad para invertir en nuevas tecnologías, mantener sistemas de monitoreo, capacitar personal y asumir los costos operativos de un sistema de monitoreo del agua.


El diagnóstico debe incluir también un análisis de los medios de transporte existentes en la zona de posible intervención y de influencia del proyecto, identificando todos los medios de transporte presentes y el estado de la infraestructura, como vías y caminos. En comunidades rurales del Cusco como Huayllarcocha, el acceso a vías de comunicación adecuadas es crucial para la implementación de proyectos, ya que determina la facilidad de transporte de equipos, materiales, personal especializado y la capacidad de mantener el sistema a través del tiempo con repuestos y mantenimiento adecuado.


La existencia de infraestructura y servicios de Salud, vivienda y educación que se prestan en la zona de posible intervención por el proyecto, en el municipio y en la posible área de influencia, debe ser documentada detalladamente. Estos servicios complementarios son importantes porque el monitoreo del agua tiene impactos directos en la salud pública (prevención de enfermedades relacionadas con agua de calidad inadecuada), en la calidad de vida de los habitantes (acceso a agua segura para consumo doméstico) y en el ámbito educativo (possibilidad de usar el sistema como herramienta de enseñanza sobre recursos hídricos).

El diagnóstico integral debe reconocer y contener información de la localidad sobre los usos, costumbres y tradiciones acerca de las formas de saneamiento básico de la población. En las comunidades andinas del Cusco, existen prácticas tradicionales de manejo del agua que han sido transmitidas generacionalmente y que deben ser comprendidas y respetadas en el diseño de cualquier intervención. La población de Huayllarcocha probablemente tiene conocimientos ancestrales sobre captación, almacenamiento y uso del agua que pueden ser integrados con las tecnologías modernas de monitoreo.



La existencia de infraestructura y delimitación de las áreas potencialmente aptas para desarrollar el proyecto de monitoreo del agua es otro elemento crítico del diagnóstico. Esto incluye identificar los sitios físicos generales donde se pueden instalar los sensores, dispositivos de monitoreo, sistemas de transmisión de datos y otros componentes del sistema. La ubicación de estos elementos debe considerar factores como la accesibilidad, la protección contra condiciones climáticas adversas, la proximidad a las fuentes de agua que se desea monitorear y la capacidad de la comunidad para acceder y mantener los equipos.

El diagnóstico debe realizar un análisis y confrontación de las demandas actuales y futuras del sistema de agua, conforme con las herramientas de planeación municipal como el POT (Plan de Ordenamiento Territorial), el PGIRS (Plan de Gestión Integral de Residuos Sólidos) y demás herramientas de planeación municipal. Este análisis permite estimar la capacidad necesaria de las obras por construirse y las expansiones futuras que requerirá el sistema de monitoreo, asegurando que el proyecto sea sostenible y escalable a largo plazo.

Una parte fundamental del diagnóstico es la identificación de los principales problemas que requieren atención en la comunidad de Huayllarcocha relacionados con el agua. Estos problemas pueden incluir falta de acceso a agua de calidad, ausencia de sistemas de monitoreo que permitan conocer el estado del recurso hídrico, enfermedades relacionadas con consumo de agua inadecuada, pérdidas en el sistema de distribución por falta de detectación temprana, y otros desafíos que limitan el desarrollo de la comunidad.

El diagnóstico participativo es esencial para incorporar la pluralidad de puntos de vista de los diferentes actores de la comunidad. Esta participación da viabilidad política, social y cultural al proyecto, brinda elementos para la jerarquización de conflictos y permite el establecimiento de prioridades, así como la definición y adopción de alternativas de solución. En Huayllarcocha, es fundamental involucrar a los habitantes locales, líderes comunitarios, autoridades locales, instituciones educativas y otros actores relevantes en el proceso de diagnóstico para asegurar que el proyecto responda a las necesidades reales de la comunidad.


El análisis del contexto social, económico y medioambiental realizado a nivel del territorio propuesto debe ser exhaustivo y detallado. Esto incluye comprender las dinámicas sociales de la comunidad, las actividades económicas principales (probablemente agricultura, ganadería y otras actividades rurales), las características medioambientales específicas de la zona (clima, topografía, recursos naturales) y cómo estos elementos interactúan con el problema de monitoreo del agua.


La metodología de diagnóstico debe seguir un enfoque estructurado que seis pasos recomendados: Paso 1. Análisis inicial para el cierre de brechas; Paso 2. Lectura sectorial y transversal del territorio; Paso 3. Identificación de problemas; Paso 4. Definición de causas y consecuencias; Paso 5. Ejercicio participativo con la comunidad; y Paso 6. Síntesis de la situación actual del territorio. Este enfoque garantiza que el diagnóstico sea completo, sistemático y participativo.


El diagnóstico debe recopilar y analizar información relevante sobre el estado actual de la entidad territorial, describiendo en detalle la situación de Huayllarcocha. Esta información debe ser documentada de manera que permita fundamentar claramente la necesidad de la intervención propuesta y justificar la inversión pública en el sistema de monitoreo del agua.

La identificación y definición de problemas es crucial para posteriormente plantear soluciones adecuadas y coherentes con los objetivos del proyecto. El problema identificado ha de ser concreto, abordable, importante para los actores implicados y necesario para el conjunto de la sociedad. En el caso de Huayllarcocha, el problema central relacionado con la ausencia de monitoreo del agua debe ser definido con precisión, analizando sus efectos (qué consecuencias provoca este problema) y sus causas (por qué ocurre ese problema).

El diagnóstico debe generar una línea base de intervención que permita medir el impacto del proyecto a lo largo del tiempo. Esta línea base incluye indicadores cuantificables sobre la situación actual del servicio de agua, la calidad del agua, la cobertura del servicio, la salud de la población relacionada con el agua, y otros elementos relevantes que puedan ser comparados con la situación futura después de la implementación del sistema de monitoreo.

La recopilarción de información para el diagnóstico debe incluir datos estadísticos actualizados, con la fecha en la cual fueron desarrollados, y las referencias utilizadas [space-instructions]. Esto es fundamental para garantizar que el diagnóstico sea basado en información factual y verificable, cumpliendo con los estándares de calidad requeridos para proyectos de inversión pública.

En resumen, el diagnóstico integral de Huayllarcocha debe proporcionar una comprensión profunda y detallada de la situación actual del territorio, la población y el estado de la infraestructura y servicios existentes, fundamentando claramente la necesidad de la intervención propuesta en el proyecto de monitoreo del agua. Este diagnóstico debe ser participativo, sistemático, exhaustivo y basado en información actualizada y verificable, cumpliendo con las normativas del Gobierno Regional del Cusco para el desarrollo de proyectos de inversión pública en el ámbito rural.


\section{Beneficiarios}

La identificación y cuantificación precisa de la población objetivo que recibirá directamente las mejoras del proyecto de monitoreo del agua en Huayllarcocha es un componente crítico para el diseño, implementación y evaluación del proyecto de inversión pública. Los beneficiarios de este proyecto se dividen en dos categorías principales: beneficiarios directos y beneficiarios indirectos, y es fundamental identificar y cuantificar cada grupo con precisión para garantizar que el proyecto responda a las necesidades reales de la comunidad y para justificar la inversión pública propuesta.

Los beneficiarios directos del proyecto de monitoreo del agua en Huayllarcocha son principalmente los habitantes locales de la comunidad que dependen directamente del agua para sus actividades domésticas, agrícolas y de subsistencia. Estos habitantes incluyen a las familias que residen permanentemente en Huayllarcocha, que utilizan el agua para consumo humano, preparación de alimentos, limpieza, higiene personal, y otras actividades esenciales del hogar. La población afectada por la carencia de monitoreo del agua será la base para determinar los beneficiarios directos del proyecto. En comunidades rurales del Cusco como Huayllarcocha, cada familia suele tener entre 4 y 6 miembros, y el número total de familias puede variar significativamente según la tamaño de la comunidad.

Además de los habitantes locales, los estudiantes vinculados a instituciones educativas de la comunidad también constituyen beneficiarios directos importantes del proyecto. Las instituciones educativas en Huayllarcocha, que pueden incluir escuelas primarias y secundarias, tienen estudiantes que dependen del agua para sus actividades educativas, consumo durante el día escolar, higiene en las instalaciones educativas y otras necesidades relacionadas con el ambiente escolar. La identificación de beneficiarios debe considerar tanto a los estudiantes como a los habitantes locales que recibirá directamente las mejoras del proyecto, según se especifica en los requisitos del proyecto [query].

La cuantificación precisa de la población beneficiaria requiere realizar un análisis detallado de la población actual de Huayllarcocha, incluyendo sus características socioeconómicas y demográficas. Esto implica recopilar datos sobre el número total de habitantes, el número de familias, la distribución por edad, el género, la estructura familiar y otros elementos demográficos que permitan caracterizar adecuadamente la población objetivo. En el contexto rural del Cusco, es importante considerar que la población puede tener características específicas como migración temporal, familias extensas que viven bajo un mismo techo, y otras dinámicas demográficas propias de las comunidades andinas.

La identificación de la población potencial es el primer paso en la cuantificación de beneficiarios. La población potencial incluye el total de la población que presenta la necesidad y/o problema que justifica el programa y pudiera ser elegible para su atención. En el caso del proyecto de monitoreo del agua, la población potencial comprende todos los habitantes de Huayllarcocha que dependen del agua para sus actividades y que actualmente no tienen acceso a un sistema de monitoreo que permita conocer el estado del recurso hídrico. Esta población potencial debe ser documentada en su totalidad antes de determinar la población objetivo específica.

La identificación de la población objetivo es el segundo paso y representa la población que el proyecto atenderá específicamente. La población objetivo puede ser igual a la población potencial si el proyecto tiene capacidad para atender a toda la población affected, o puede ser un subgrupo de la población potencial si el proyecto tiene limitaciones de recursos, cobertura o capacidad operativa. En el caso de Huayllarcocha, es probable que la población objetivo incluya a toda la población de la comunidad, dado que un sistema de monitoreo del agua puede beneficiar a toda la comunidad simultáneamente sin limitaciones de cobertura significativas.

Para cuantificar precisamente los beneficiarios directos, es necesario considerar varios elementos. Primero, el número total de habitantes de Huayllarcocha debe ser determinado mediante datos actualizados del Instituto Nacional de Estadística y Informática (INEI) del Perú, o mediante censos comunitarios realizados recientemente. Segundo, el número de estudiantes en las instituciones educativas de la comunidad debe ser identificado, incluyendo tanto estudiantes de nivel primario como secundario. Tercero, es importante considerar si hay población temporal o visitante que también pueda beneficiarse del sistema de monitoreo, aunque estos generalmente se consideran beneficiarios indirectos.

Las características socioeconómicas de la población beneficiaria deben ser detalladamente documentadas, incluyendo la población actual, estratificación socioeconómica, índice de NBI (Necesidad Básica Insatisfecha), población en miseria, y otros indicadores de bienestar. Estas características son importantes porque determinan la vulnerabilidad de la población frente a problemas relacionados con la calidad del agua, la capacidad de la comunidad para invertir en mantenimiento del sistema, y la prioridad que debe tener el proyecto en términos de impacto social.

En las comunidades rurales del Cusco, la población suele tener índices de NBI significativos, lo que indica que muchas familias no satisface necesidades básicas como acceso a agua segura, servicios de salud, educación de calidad y vivienda adecuada. Esta característica de la población beneficiaria justifica la inversión pública en el proyecto de monitoreo del agua, dado que el proyecto contribuye directamente al cierre de brechas en necesidades básicas insatisfechas, especialmente en el acceso a agua de calidad.

La distribución por edad de la población beneficiaria es otro elemento importante de la cuantificación. En comunidades andinas como Huayllarcocha, es común tener una población con proporción significativa de niños y adultos jóvenes, debido a tasas de natalidad relativamente altas y migración de adultos jóvenes hacia centros urbanos. Esta distribución por edad es importante porque los niños son particularmente vulnerables a enfermedades relacionadas con agua de calidad inadecuada, y los adultos jóvenes son los que probablemente operarán y mantendrán el sistema de monitoreo a lo largo del tiempo.

El género de la población beneficiaria también debe ser considerado en la cuantificación. En las comunidades rurales del Cusco, las mujeres suelen tener roles principales en el manejo del agua para actividades domésticas, preparación de alimentos, limpieza y cuidado de la familia. Por lo tanto, las mujeres son beneficiarias directas particularly importantes del proyecto de monitoreo del agua, dado que el sistema les permitirá conocer la calidad del agua que utilizan para estas actividades esenciales y tomar decisiones informadas sobre su uso.

La cuantificación de beneficiarios debe incluir también una estimación de los beneficiarios indirectos del proyecto. Los beneficiarios indirectos son personas que no reciben directamente las mejoras del proyecto pero que experimentan efectos positivos derivados de la implementación del sistema de monitoreo. En el caso de Huayllarcocha, los beneficiarios indirectos pueden incluir a poblaciones de comunidades vecinas que puedan acceder a la información generada por el sistema de monitoreo, autoridades locales y regionales que utilizan la información para tomar decisiones sobre gestión de recursos hídricos, y organismos de salud que monitorean enfermedades relacionadas con agua.

La identificación de beneficiarios debe considerar también la área de influencia del proyecto, que incluye no solo a la comunidad de Huayllarcocha sino también a poblaciones en municipios cercanos que puedan beneficiarse de la información generada por el sistema de monitoreo. Esta área de influencia puede extenderse a varias comunidades rurales en la región del Cusco que comparten fuentes de agua o que tienen interés en la gestión de recursos hídricos.

Para garantizar una cuantificación precisa, es necesario utilizar metodologías de recolección de datos que sean apropiadas para el contexto rural de Huayllarcocha. Estas metodologías pueden incluir censos comunitarios, entrevistas a líderes locales, análisis de registros educativos y consulta con autoridades locales. La participación de la comunidad en el proceso de cuantificación es fundamental para asegurar que los datos sean precisos y que la comunidad se sienta involucrada en el proyecto desde etapas tempranas.

La cuantificación de beneficiarios debe ser actualizada periódicamente durante la implementación del proyecto, dado que la población de comunidades rurales puede cambiar debido a factores como migración, nacimientos, muertes y otros movimientos demográficos. Es importante establecer un sistema de monitoreo y evaluación que permita actualizar regularmente el número de beneficiarios y medir el impacto del proyecto en la población objetivo.

Los beneficiarios del proyecto de monitoreo del agua en Huayllarcocha recibirán beneficios directos que incluyen: conocimiento preciso sobre la calidad del agua que consumen, capacidad para detectar problemas de contaminación o deterioro de la calidad del agua, información para tomar decisiones informadas sobre el uso del agua, prevención de enfermedades relacionadas con agua de calidad inadecuada, y mejora en la gestión de los recursos hídricos de la comunidad. Estos beneficios directos justifican la inclusión de la población identificada como beneficiaria del proyecto.

Además de los beneficios directos, los beneficiarios del proyecto recibiran beneficios indirectos que incluyen: mejora en la salud pública de la comunidad, aumento en la confianza de la población sobre la seguridad del agua, capacidad para planificar actividades agrícolas y de subsistencia basadas en información precisa sobre recursos hídricos, fortalecimiento de la capacidad organizacional de la comunidad para gestionar recursos comunes, y posibilidad de usar el sistema de monitoreo como herramienta educativa en instituciones locales.

La identificación y cuantificación de beneficiarios debe cumplir con los requisitos normativos del Gobierno Regional del Cusco para proyectos de inversión pública. Estos requisitos incluyen la documentación detallada de la población objetivo, la justificación de la elegibilidad de los beneficiarios, la demostración de que el proyecto responde a necesidades reales de la comunidad, y la evidencia de que la inversión pública tendrá un impacto positivo significativo en la población beneficiaria.

En conclusión, los beneficiarios del proyecto de monitoreo del agua en Huayllarcocha son principalmente los habitantes locales de la comunidad y los estudiantes de las instituciones educativas locales, con una cuantificación precisa que debe basarse en datos actualizados y verificables sobre la población actual de la comunidad. La identificación y cuantificación de beneficiarios es fundamental para fundamentar la inversión pública, garantizar que el proyecto responda a necesidades reales, y medir el impacto del proyecto a lo largo del tiempo.


\section{Problema Central}

La definición clara y concreta de la necesidad o carencia principal que limita el desarrollo de la comunidad de Huayllarcocha y que se pretende resolver constituye el problema central del proyecto de inversión pública de monitoreo del agua. El problema central debe ser identificado con precisión, siendo concreto, abordable, importante para los actores implicados y necesario para el conjunto de la sociedad. En el contexto de Huayllarcocha, una comunidad rural del Cusco, el problema central se identifica como la ausencia de sistemas de monitoreo del agua que permitan conocer en tiempo real la calidad, disponibilidad y estado del recurso hídrico que utiliza la comunidad para sus actividades domésticas, agrícolas y de subsistencia.

La ausencia de monitoreo del agua en Huayllarcocha representa una carencia tecnológica crítica que limita significativamente la capacidad de la comunidad para gestionar adecuadamente sus recursos hídricos. Sin sistemas de monitoreo, la población no tiene acceso a información precisa sobre la calidad del agua que consume, no puede detectar tempranamente problemas de contaminación, no conoce la disponibilidad real del recurso en diferentes periods del año, y no puede tomar decisiones basadas en datos sobre el uso del agua. Esta falta de información genera una situación de vulnerabilidad donde la comunidad depende de percepciones subjetivas, observaciones visuales limitadas y conocimientos tradicionales que, aunque valiosos, no son suficientes para garantizar el acceso a agua segura y para la gestión óptima del recurso.

El problema central de la ausencia de monitoreo del agua es concreto porque se puede identificar claramente como la falta de dispositivos, sensores, sistemas de medición y tecnologías que permitan recopilar, procesar y transmitir datos sobre el estado del agua. Es abordable porque existen soluciones tecnológicas disponibles, tanto nacionales como internacionales, que pueden ser implementadas en el contexto rural de Huayllarcocha con recursos financieros razonables y capacidad técnica accesible. Es importante para los actores implicados porque la población de Huayllarcocha, las autoridades locales, las instituciones educativas y otros stakeholders reconocen que la calidad y disponibilidad del agua es fundamental para su salud, sus actividades económicas y su desarrollo general. Es necesario para el conjunto de la sociedad porque el acceso a agua segura es un derecho humano fundamental y la gestión adecuada de recursos hídricos contribuye al desarrollo sostenible de la región del Cusco.

La carencia de monitoreo del agua en Huayllarcocha genera múltiples consecuencias negativas que limitan el desarrollo de la comunidad. Primero, la población consume agua sin conocer su calidad real, lo que aumenta el riesgo de enfermedades relacionadas con contaminación del agua, incluyendo enfermedades gastrointestinales, parasitosis, y otras patologías que afectan particularmente a niños, mujeres embarazadas y personas con sistemas inmunitarios debilitados. Segundo, la comunidad no puede detectar tempranamente problemas de contaminación del agua, lo que significa que cuando se identifican problemas de calidad, ya pueden haber ocurrido exposiciones prolongadas a agua inadecuada. Tercero, la falta de información sobre disponibilidad del agua limita la capacidad de la comunidad para planificar actividades agrícolas, de ganado y otras actividades de subsistencia que dependen del recurso hídrico.

El problema de ausencia de monitoreo del agua también limita el desarrollo de la comunidad en términos de capacidad organizacional y gestión de recursos comunes. Sin información precisa sobre el estado del agua, la comunidad no puede tomar decisiones colectivas informadas sobre distribución del agua, mantenimiento de sistemas de distribución, inversiones en mejora de infraestructura, o implementación de prácticas de conservación del recurso. Esta limitación en la toma de decisiones basada en datos fortalece dinámicas de gestión reactiva rather than preventiva, donde la comunidad responde a problemas cuando ya se han manifestado completamente, rather than anticipando y previenen problemas.

La carencia de monitoreo del agua en Huayllarcocha es particularmente crítica dado el contexto rural de la comunidad y las características específicas de la región del Cusco. En zonas rurales altoandinas, las fuentes de agua suelen ser ríos, quebradas, lagunas o sistemas de captación de agua de lluvia que pueden ser susceptibles a contaminación por actividades agrícolas, ganaderas, descargas no controladas, o cambios en las condiciones climáticas. Sin sistemas de monitoreo, la comunidad no tiene capacidad para detectar estos cambios en la calidad del agua ni para implementar medidas correctivas oportunas.

El problema central de ausencia de monitoreo del agua también se relaciona con brechas importantes en necesidades básicas insatisfechas (NBI) que caracterizan a muchas comunidades rurales del Cusco. El acceso a agua segura es un componente fundamental de las necesidades básicas, y la falta de monitoreo contribuye a mantener esta brecha porque la comunidad no puede garantizar que el agua que consume es segura. Esta situación perpetúa condiciones de vulnerabilidad social y limita el desarrollo integral de la comunidad.

La definición del problema central debe considerar también las causas que originan la ausencia de monitoreo del agua en Huayllarcocha. Estas causas incluyen: limitaciones económicas que impiden la inversión en tecnologías de monitoreo, falta de conocimiento técnico sobre sistemas de monitoreo disponibles, ausencia de capacitación en operación y mantenimiento de sistemas de monitoreo, limitaciones en el acceso a proveedores de tecnologías de monitoreo en zonas rurales, y falta de priorización de este tema en la planificación de desarrollo comunitario. La comprensión de estas causas es fundamental para diseñar soluciones que aborden no solo el problema superficial (ausencia de monitoreo) sino también las causas raíz que lo generan.

El problema central de ausencia de monitoreo del agua en Huayllarcocha representa una oportunidad importante para la inversión pública porque su resolución genera impactos positivos significativos en múltiples dimensiones del desarrollo comunitario. La implementación de un sistema de monitoreo del agua permitirá a la comunidad conocer la calidad del agua que consume, detectar tempranamente problemas de contaminación, planificar actividades basadas en información precisa sobre disponibilidad del recurso, tomar decisiones colectivas informadas sobre gestión del agua, prevenir enfermedades relacionadas con agua inadecuada, y fortalecer la capacidad organizacional de la comunidad para gestionar recursos comunes.

La claridad y concreta del problema central es fundamental para el diseño del proyecto porque permite enfocar los esfuerzos de intervención en la resolución de una necesidad específica y medible. El problema de ausencia de monitoreo del agua es claramente identificable, puede ser medido mediante indicadores cuantificables (por ejemplo, número de dispositivos de monitoreo instalados, frecuencia de medición de calidad del agua, porcentaje de población con acceso a información sobre calidad del agua), y su resolución puede ser evaluada de manera objetiva mediante indicadores de impacto.

En el contexto de la normativa del Gobierno Regional del Cusco para proyectos de inversión pública en el ámbito rural, el problema central de ausencia de monitoreo del agua en Huayllarcocha cumple con los requisitos de justificación de inversión pública porque representa una necesidad real de la comunidad, tiene impacto significativo en la salud y desarrollo de la población, contribuye al cierre de brechas en necesidades básicas insatisfechas, y puede ser resuelto mediante intervención pública con recursos financieros razonables y capacidad técnica disponible.


\section{Árbol de Causas y Efectos}

El árbol de causas y efectos es un diagrama estructurado que visualiza las raíces que originan el problema central de ausencia de monitoreo del agua en Huayllarcocha y las consecuencias negativas que este problema genera en la población. Este herramienta metodológica permite comprender la relación causal entre el problema central, sus causas fundamentais y sus efectos negativos, facilitando el diseño de soluciones que aborden tanto el problema superficial como las causas raíz.

\subsection{Causas del problema central}
La causa fundamental número uno es la limitación económica de la comunidad. Huayllarcocha, como comunidad rural del Cusco, presenta características socioeconómicas que incluyen índice de NBI significativo, estratificación socioeconómica baja y condiciones de pobreza que limitan el acceso a tecnologías avanzadas. Los habitantes de la comunidad tienen ingresos limitados derivados principalmente de actividades agrícolas y de subsistencia, lo que significa que no tienen capacidad financiera para invertir en sistemas de monitoreo del agua que requieren dispositivos, sensores, equipos de transmisión de datos y otros componentes tecnológicos con costos de adquisición y mantenimiento.

La causa fundamental número dos es la falta de conocimiento técnico sobre sistemas de monitoreo disponibles. En comunidades rurales como Huayllarcocha, la población tiene conocimientos tradicionales sobre manejo del agua transmitidos generacionalmente, pero carece de conocimiento sobre tecnologías modernas de monitoreo disponibles en el mercado nacional e internacional. Esta falta de conocimiento técnico incluye no saber qué sistemas de monitoreo existen, cómo funcionan, qué costos tienen, qué beneficios generan, cómo se operan y mantienen, y qué proveedores existen. Esta brecha de conocimiento limita la capacidad de la comunidad para identificar, evaluar e implementar soluciones tecnológicas apropiadas.

La causa fundamental número tres es la ausencia de capacitación en operación y mantenimiento. Los sistemas de monitoreo del agua requieren personal capacitado para operar los dispositivos, interpretar los datos generados, realizar mantenimiento preventivo y correctivo, y tomar decisiones basadas en la información obtenida. En Huayllarcocha, no existen programas de capacitación formal o informal que preparen a la población para estas tareas, lo que significa que incluso si la comunidad adquiere sistemas de monitoreo, no tendría capacidad para operarlos y mantenerlos adecuadamente a lo largo del tiempo.

La causa fundamental número cuatro es las limitaciones en el acceso a proveedores de tecnologías en zonas rurales. Las comunidades rurales del Cusco como Huayllarcocha tienen acceso limitado a proveedores de tecnologías de monitoreo debido a su ubicación geográfica, distancia a centros urbanos, y falta de infraestructura de transporte adecuada. Los proveedores de sistemas de monitoreo generalmente se concentran en ciudades principales como Cusco, y el acceso a estos proveedores requiere transporte, tiempo y recursos que la comunidad rural no tiene fácilmente disponibles. Esta limitación en el acceso a proveedores dificulta la adquisición de equipos, la obtención de repuestos, y el acceso a servicios de soporte técnico.

La causa fundamental número cinco es la falta de priorización en la planificación de desarrollo comunitario. En comunidades rurales como Huayllarcocha, la planificación de desarrollo suele priorizar necesidades inmediatas como alimentación, vivienda, salud básica y educación, dejando temas como monitoreo de recursos hídricos como prioritarios secundarios o tertiary. Esta falta de priorización significa que no se asignan recursos financieros, tiempo ni esfuerzo institucional para implementar sistemas de monitoreo del agua, incluso cuando la comunidad reconoce que este es un tema importante.

\subsection{Efectivos negativos del problema central}

El efecto negativo número uno es el aumento del riesgo de enfermedades relacionadas con contaminación del agua. Cuando la comunidad consume agua sin conocer su calidad real, aumenta significativamente el riesgo de enfermedades gastrointestinales, parasitosis, y otras patologías relacionadas con agua de calidad inadecuada. Estos enfermedades afectan particularmente a niños, mujeres embarazadas y personas con sistemas inmunitarios debilitados, generando impactos negativos en la salud pública de la comunidad, aumento en costos de atención médica, reducción en la capacidad de trabajo de los adultos, y disminución en el rendimiento escolar de los niños.

El efecto negativo número dos es la detección tardía de problemas de contaminación del agua. Sin sistemas de monitoreo que permitan detectar tempranamente problemas de contaminación, la comunidad solo identifica problemas de calidad del agua cuando ya se han manifestado completamente, generalmente cuando aparecen síntomas de enfermedades en la población o cuando el agua presenta cambios visibles obvios. Esta detección tardía significa que la comunidad puede haber consumido agua contaminada durante periodos prolongados antes de identificar el problema, aumentando exponencialmente el riesgo de exposición y enfermedades.

El efecto negativo número tres es la limitación en la planificación de actividades agrícolas y de subsistencia. La falta de información sobre disponibilidad del agua en diferentes periods del año limita significativamente la capacidad de la comunidad para planificar actividades agrícolas, de ganado y otras actividades de subsistencia que dependen del recurso hídrico. Los agricultores no pueden conocer cuándo habrá suficiente agua para cultivos, no pueden planificar ciclos de cultivo basados en disponibilidad real del agua, no pueden tomar decisiones óptimas sobre qué cultivos sembrar, y no pueden gestionar eficientemente el recurso hídrico disponible.

El efecto negativo número cuatro es la tomada de decisiones colectivas no informadas sobre gestión del agua. Sin información precisa sobre el estado del agua, la comunidad no puede tomar decisiones colectivas informadas sobre distribución del agua, mantenimiento de sistemas de distribución, inversiones en mejora de infraestructura, o implementación de prácticas de conservación del recurso. Esto genera dinámicas de gestión reactiva rather than preventiva, conflictos comunitarios sobre distribución del agua, uso inefficiente del recurso, y falta de coordinación en acciones de conservación.

El efecto negativo número cinco es la perpetuación de brechas en necesidades básicas insatisfechas (NBI). La falta de monitoreo contribuye a mantener brechas en acceso a agua segura, que es un componente fundamental de las necesidades básicas. Esta situación perpetúa condiciones de vulnerabilidad social en la comunidad, limita el desarrollo integral de Huayllarcocha, y contribuye a mantener la comunidad en condiciones de pobreza y NBI significativo.

El efecto negativo número seis es la debilidad en la capacidad organizacional para gestionar recursos comunes. La ausencia de información precisa sobre el estado del agua limita el fortalecimiento de la capacidad organizacional de la comunidad para gestionar recursos hídricos como activo común. Sin datos concretos, la comunidad no puede desarrollar mecanismos de gestión colectiva efectivos, no puede establecer normas claras sobre uso del agua, no puede implementar sistemas de monitoreo comunitario, y no puede fortalecer su organización para la gestión de recursos naturales.

El efecto negativo número siete es la vulnerabilidad frente a cambios climáticos y ambientales. Las comunidades rurales del Cusco son particularmente vulnerables a cambios climáticos que afectan la disponibilidad y calidad del agua, incluyendo cambios en patrones de lluvia, aumento en temperatura, y eventos climáticos extremos. Sin sistemas de monitoreo, la comunidad no puede detectar estos cambios tempranamente, no puede adaptar sus prácticas de uso del agua, y no puede implementar medidas de mitigación o adaptación oportunas.

\subsection{Relación Causal}

La relación causal en el árbol de causas y efectos muestra cómo las causas fundamentales generan el problema central de ausencia de monitoreo del agua, y cómo este problema central genera múltiples efectos negativos en la población de Huayllarcocha. Las causas económicas, técnicas, de capacitación, de acceso a proveedores y de priorización se combinan para crear la situación de ausencia de monitoreo. Este problema central, a su vez, genera efectos en salud pública, detección de contaminación, planificación agrícola, toma de decisiones, perpetuación de NBI, capacidad organizacional y vulnerabilidad climática.

La comprensión de esta relación causal es fundamental para diseñar soluciones que aborden tanto el problema superficial (ausencia de monitoreo) como las causas raíz (limitaciones económicas, falta de conocimiento, ausencia de capacitación, limitaciones de acceso, falta de priorización). Soluciones que solo abordan el problema superficial sin atender las causas tendrán limitada efectividad y sostenibilidad a largo plazo.

\section{Objetivo General}

El objetivo general del proyecto de monitoreo del agua en Huayllarcocha es el enunciado que describe el cambio positivo principal que el proyecto logrará al solucionar el problema central identificado de ausencia de sistemas de monitoreo del agua. El objetivo general debe ser claro, específico, medible, alcanzable, relevante y temporalmente definido, cumpliendo con los estándares de calidad requeridos para proyectos de inversión pública bajo la normativa del Gobierno Regional del Cusco.

El objetivo general del proyecto es: "Implementar un sistema de monitoreo del agua en la comunidad de Huayllarcocha que permita conocer en tiempo real la calidad, disponibilidad y estado del recurso hídrico, facilitando la toma de decisiones informadas sobre gestión del agua, prevenciendo enfermedades relacionadas con agua de calidad inadecuada, y fortaleciendo la capacidad de la comunidad para gestionar adecuadamente sus recursos hídricos de manera sostenible."

Este objetivo general describe el cambio positivo principal que el proyecto logrará: la implementación de un sistema de monitoreo funcional que transforma la situación actual de ausencia de información sobre el agua hacia una situación de conocimiento preciso y en tiempo real sobre calidad, disponibilidad y estado del recurso hídrico. El objetivo es claro porque especifica exactamente qué se implementará (sistema de monitoreo del agua), dónde (comunidad de Huayllarcocha), y qué información se generará (calidad, disponibilidad y estado del recurso hídrico en tiempo real).

El objetivo general es específico porque identifica los componentes principales del cambio positivo: implementación de sistema de monitoreo, conocimiento en tiempo real, toma de decisiones informadas, prevención de enfermedades, y fortalecimiento de capacidad comunitaria. Cada componente es identificable y puede ser evaluado de manera objetiva mediante indicadores cuantificables.

El objetivo general es medible porque puede ser evaluado mediante indicadores como: número de dispositivos de monitoreo instalados y funcionando, frecuencia de medición de calidad del agua (por ejemplo, mediciones diarias, semanales o mensuales), porcentaje de población con acceso a información sobre calidad del agua, número de enfermedades relacionadas con agua que disminuyen después de la implementación del sistema, y nivel de capacidad organizacional de la comunidad para gestionar recursos hídricos (evaluado mediante indicadores de organización comunitaria).

El objetivo general es alcanzable porque existen soluciones tecnológicas disponibles, tanto nacionales como internacionales, que pueden ser implementadas en el contexto rural de Huayllarcocha con recursos financieros razonables y capacidad técnica accesible. Sistemas de monitoreo del agua con sensores, dispositivos de medición automática y sistemas de transmisión de datos en tiempo real están disponibles en el mercado y pueden ser adaptados al contexto rural de comunidades andinas del Cusco.

El objetivo general es relevante porque responde directamente al problema central identificado (ausencia de monitoreo del agua) y genera impactos positivos significativos en múltiples dimensiones del desarrollo comunitario: salud pública, planificación de actividades agrícolas, toma de decisiones colectivas, cierre de brechas en NBI, fortalecimiento de capacidad organizacional, y reducción de vulnerabilidad frente a cambios climáticos.

El objetivo general es temporalmente definido porque puede ser implementado dentro de un periodo razonable de tiempo (generalmente 12-24 meses para proyectos de inversión pública de esta naturaleza), con fases claras de diseño, implementación, capacitación y evaluación.

El objetivo general del proyecto contribuye directamente al desarrollo de la comunidad de Huayllarcocha porque aborda una necesidad fundamental (acceso a información sobre calidad del agua) que tiene impactos positivos en múltiples dimensiones del bienestar comunitario. La implementación del sistema de monitoreo permitirá a la comunidad conocer la calidad del agua que consume, lo que contribuirá a la prevención de enfermedades relacionadas con agua inadecuada, mejorando la salud pública general de la comunidad.

El objetivo general también contribuye al cierre de brechas en necesidades básicas insatisfechas (NBI) que caracterizan a muchas comunidades rurales del Cusco. El acceso a agua segura y la capacidad de conocer su calidad son componentes fundamentales de las necesidades básicas, y la implementación del sistema de monitoreo contribuye directamente a cerrar estas brechas.

La implementación del sistema de monitoreo según el objetivo general fortalecerá la capacidad organizacional de la comunidad para gestionar recursos comunes, un elemento crucial para el desarrollo sostenible de Huayllarcocha. Con información precisa sobre el estado del agua, la comunidad podrá tomar decisiones colectivas informadas, establecer normas claras sobre uso del agua, implementar mecanismos de gestión colectiva efectivos, y fortalecer su organización para la gestión de recursos naturales.

El objetivo general del proyecto es consistente con la normativa del Gobierno Regional del Cusco para proyectos de inversión pública en el ámbito rural porque representa una necesidad real de la comunidad, tiene impacto significativo en la salud y desarrollo de la población, contribuye al desarrollo sostenible de la región, y puede ser resuelto mediante intervención pública con recursos financieros razonables y capacidad técnica disponible.


\section{Objetivos Especificos}

Los objetivos específicos son metas operativas y medibles relacionadas con infraestructura, equipamiento y capacitación necesarias para alcanzar el propósito general del proyecto de monitoreo del agua en Huayllarcocha. Estos objetivos específicos desglosan el objetivo general en componentes más concretos y accionables, cada uno con indicadores cuantificables que permiten evaluar el progreso y el éxito del proyecto.


\begin{itemize}
    \item Infraestructura
    \item Equipamiento
    \item Capacitación
    \item Sistema de información
    \item Sostenibilidad
\end{itemize}

\subsection{Infraestructura}
"Construir y instalar la infraestructura física necesaria para el sistema de monitoreo del agua, incluyendo puntos de muestreo, estaciones de monitoreo, protección contra condiciones climáticas, y sistemas de energía, asegurando que la infraestructura sea durable, segura y accesible para la comunidad de Huayllarcocha."

Este objetivo específico se enfoca en la componente de infraestructura del proyecto. Los indicadores medibles incluyen: número de puntos de muestreo construidos (por ejemplo, 3-5 puntos en diferentes fuentes de agua de la comunidad), número de estaciones de monitoreo instaladas (por ejemplo, 2-3 estaciones principales), calidad de la infraestructura construida (evaluada mediante estándares de durabilidad y seguridad), y accesibilidad de la infraestructura para la comunidad (evaluada mediante distancia a puntos de muestreo, facilidad de acceso, y seguridad de los instalaciones).

La infraestructura física incluye puntos de muestreo donde se instalaran los sensores y dispositivos de monitoreo, estaciones de monitoreo que protegen los equipos de condiciones climáticas adversas (lluvia, nieve, viento, temperatura extrema), sistemas de protección contra condiciones climáticas que garantizan la operatividad continua del sistema, y sistemas de energía (paneles solares, baterías, o conexiones eléctricas) que proporcionan energía constante para el funcionamiento de los dispositivos de monitoreo.

\subsection{Equipamiento}

"Adquirir y instalar el equipamiento tecnológico necesario para el sistema de monitoreo del agua, incluyendo sensores de calidad del agua, dispositivos de medición automática, sistemas de transmisión de datos en tiempo real, equipos de procesamiento de información, y dispositivos de Visualización de datos, asegurando que el equipamiento sea apropiado para el contexto rural de Huayllarcocha, preciso, durable y de costo razonable."

Este objetivo específico se enfoca en la componente de equipamiento del proyecto. Los indicadores medibles incluyen: número de sensores de calidad del agua adquiridos e instalados (por ejemplo, sensores para pH, temperatura, turbiedad, oxígeno disuelto, conductividad), número de dispositivos de medición automática instalados (por ejemplo, 2-3 dispositivos principales), capacidad de transmisión de datos en tiempo real (evaluada mediante frecuencia de transmisión y confiabilidad de la conexión), precisión de las mediciones (evaluada mediante comparación con estándares de referencia), y durabilidad del equipamiento (evaluada mediante periodo de funcionamiento continuo sin fallas).

El equipamiento tecnológico incluye sensores de calidad del agua que miden parámetros críticos como pH, temperatura, turbiedad, oxígeno disuelto, conductividad, y otros indicadores de calidad; dispositivos de medición automática que recopilan datos de los sensores de manera continua y automática; sistemas de transmisión de datos en tiempo real que transmiten la información recopilada a plataformas de Visualización accesibles para la comunidad; equipos de procesamiento de información que analizan los datos recopilados y generan informes y alertas; y dispositivos de Visualización de datos como pantallas, tablets o aplicaciones móviles que permiten a la comunidad acceder a la información sobre calidad del agua.

\subsection{Capacitación}

"Implementar un programa de capacitación completo para la población de Huayllarcocha que incluya capacitación en operación del sistema de monitoreo, interpretación de datos de calidad del agua, mantenimiento preventivo y correctivo de equipos, toma de decisiones basadas en información, y gestión comunitaria de recursos hídricos, asegurando que al menos 80% de la población objetivo haya recibido capacitación adecuada y tenga capacidad para operar y mantener el sistema de monitoreo."

Este objetivo específico se enfoca en la componente de capacitación del proyecto. Los indicadores medibles incluyen: porcentaje de población objetivo que ha recibido capacitación (objetivo: 80% o más), número de sesiones de capacitación realizadas (por ejemplo, 10-15 sesiones distribuidas en 12 meses), contenido de la capacitación cubierto (evaluado mediante checklist de temas capacitados), nivel de comprensión de los capacitados (evaluado mediante pruebas de conocimiento antes y después de la capacitación), y capacidad operativa de los capacitados (evaluada mediante demostración práctica de operación y mantenimiento del sistema).

El programa de capacitación incluye capacitación en operación del sistema de monitoreo que enseña a la población cómo activar, usar y interactuar con los dispositivos de monitoreo; capacitación en interpretación de datos de calidad del agua que enseña a comprender los parámetros medidos, interpretar valores normales y anomalías, y entender qué significan los diferentes niveles de calidad del agua; capacitación en mantenimiento preventivo y correctivo de equipos que enseña a realizar limpieza regular, inspección periódica, reemplazo de componentes cuando sea necesario, y reparación de fallas menores; capacitación en toma de decisiones basadas en información que enseña a utilizar la información generada por el sistema para tomar decisiones sobre uso del agua, distribución, conservación y otras acciones relacionadas; y capacitación en gestión comunitaria de recursos hídricos que enseña a desarrollar mecanismos de gestión colectiva, establecer normas sobre uso del agua, y fortalecer la organización comunitaria para la gestión de recursos naturales.


\subsection{Sistema de Información}

"Desarrollar y implementar un sistema de información accesible para la comunidad de Huayllarcocha que incluya plataforma digital para Visualización de datos en tiempo real, generación de informes periódicos sobre calidad del agua, sistema de alertas tempranas para problemas de contaminación, y documentación técnica del sistema de monitoreo, asegurando que el sistema de información sea fácil de usar, accesible para toda la población, y actualizado regularmente."

Este objetivo específico se enfoca en la creación de un sistema de información que permita a la comunidad acceder y utilizar la información generada por el sistema de monitoreo. Los indicadores medibles incluyen: funcionalidad de la plataforma digital (evaluada mediante disponibilidad continua y accesibilidad), frecuencia de actualización de datos (objetivo: actualización en tiempo real o diaria), calidad de los informes periódicos generados (evaluada mediante claridad, utilidad y relevancia de la información), efectividad del sistema de alertas tempranas (evaluada mediante tiempo de respuesta a problemas de contaminación), y disponibilidad de documentación técnica completa (evaluada mediante completeness de la documentación y accesibilidad para la comunidad).

El sistema de información incluye plataforma digital para Visualización de datos en tiempo real que permite a la comunidad acceder a información actualizada sobre calidad del agua mediante computador, tablet o dispositivo móvil; generación de informes periódicos sobre calidad del agua que proporcionan análisis resumidos de los datos recopilados, tendencias observadas, y recomendaciones para la comunidad; sistema de alertas tempranas para problemas de contaminación que notifica automáticamente a la comunidad cuando se detectan niveles de calidad del agua fuera de rangos normales; y documentación técnica del sistema de monitoreo que incluye manuales de operación, procedimientos de mantenimiento, especificaciones técnicas de equipos, y guías de uso para la comunidad.

\subsection{Sostenibilidad}
"Establecer mecanismos de sostenibilidad financeira, técnica y organizacional para el sistema de monitoreo del agua que incluyan plan de financiamiento para mantenimiento y operación a largo plazo, sistema de gestión de repuestos y mantenimiento, organización comunitaria responsable del sistema, y enlace con instituciones locales y regionales para soporte continuo, asegurando que el sistema de monitoreo sea sostenible a lo largo de al menos 5 años después de la implementación inicial."

Este objetivo específico se enfoca en garantizar la sostenibilidad del proyecto a lo largo del tiempo. Los indicadores medibles incluyen: existencia de plan de financiamiento para mantenimiento y operación a largo plazo (evaluada mediante completeness y viabilidad del plan), sistema de gestión de repuestos y mantenimiento establecido (evaluada mediante disponibilidad de repuestos y procedimientos de mantenimiento), organización comunitaria responsable del sistema formalizada (evaluada mediante existencia de estructura organizacional clara y roles definidos), enlace con instituciones locales y regionales establecido (evaluada mediante number de instituciones vinculadas y frecuencia de interacción), y funcionamiento continuo del sistema después de 5 años (evaluada mediante operación continua sin fallas significativas).

Los mecanismos de sostenibilidad incluyen plan de financiamiento para mantenimiento y operación a largo plazo que identifica fuentes de recursos financieras para costos operativos continuos (puede incluir contribuciones comunitarias, apoyo de autoridades locales, programas gubernamentales, o otros mecanismos); sistema de gestión de repuestos y mantenimiento que establece procedimientos para adquisición de repuestos, realización de mantenimiento regular, y respuesta a fallas; organización comunitaria responsable del sistema que formaliza una estructura organizacional con roles, responsabilidades y mecanismos de toma de decisiones para la gestión del sistema; y enlace con instituciones locales y regionales para soporte continuo que establece relaciones con autoridades locales, gobierno regional, instituciones educativas, y otras organizaciones que pueden proporcionar apoyo técnico, financiero o institucional a lo largo del tiempo.


\section{Alternativas de solución}

Las alternativas de solución son diversas opciones tecnológicas para el sistema de monitoreo del agua que deben ser comparadas técnica y económicamente para seleccionar la propuesta más viable y eficiente para el contexto rural de Huayllarcocha, considerando las limitantes establecidas por el Gobierno Regional del Cusco para el desarrollo de proyectos en el ámbito rural.

Alternativa 1: Sistema de Monitoreo Convencional con Sensores Analógicos

Esta alternativa utiliza sensores analógicos convencionales que requieren medición manual periódica por personal de la comunidad. Los sensores miden parámetros como pH, temperatura, turbiedad y otros indicadores de calidad del agua, pero los datos deben ser recopilados manualmente mediante lectura de instrumentos y registro en formatos físicos. El sistema no tiene transmisión automática de datos ni capacidad de monitoreo en tiempo real.

\subsection{Características técnicas}
\begin{itemize}
    \item Sensores analógicos de pH, temperatura, turbiedad
    \item Equipos de medición manual (medidores portátiles)
    \item Formularios físicos para registro de datos
    \item Sin transmisión automática de datos
    \item Monitoreo periódico manual (por ejemplo, semanal o mensual)
    \item Sin capacidad de alertas automáticas
    \item Requiere presencia humana constante para operación
\end{itemize}

\subsection{Viabilidad para Huayllarcocha}
\begin{itemize}
    \item Adecuación al contexto rural: Alta - tecnología simple, fácil de entender y operar
    \item Requisitos de capacitación: Baja - capacitación mínima requerida
    \item Accesibilidad de proveedores: Alta - equipos analógicos disponibles ampliamente
    \item Mantenimiento: Baja complejidad - mantenimiento simple, reemplazo de sensores frecuente
    \item Sostenibilidad: Media - depende de presencia continua de personal para medición manual
\end{itemize}


\subsection{Limitantes}
\begin{itemize}
    \item No proporciona monitoreo en tiempo real
    \item Requiere presencia humana constante para operación
    \item Datos no disponibles inmediatamente para la comunidad
    \item Sin capacidad de alertas automáticas
    \item Posible error humano en medición y registroDetección tardía de problemas de contaminación
\end{itemize}












